\section{An independent formalization of the positive-shift family in A211417}
\label{sec:a211417-positive-shifts}

\subsection{Statement and scope}

For $n\ge0$, let
\[
 A(n)=\frac{(30n)!\,n!}{(15n)!\,(10n)!\,(6n)!}.
\]
This is the classical integral factorial ratio recorded as OEIS A211417.
In a comment dated August 28, 2025, Peter Bala conjectured that for each
$k\in\{2,3,5\}$ and each integer $r\ge1$, there is a positive constant
$C(k,r)$ such that
\[
 \frac{C(k,r)A(n)}
 {\displaystyle\prod_{\substack{1\le i\le r\\\gcd(i,k)=1}}(kn+i)}
 \quad\hbox{is an integer for every }n\ge0.
\]
The checked primary source is A211417, revision 94, dated August 26, 2026.
The entire positive-shift assertion was already proved in the September
19, 2026 report \emph{Polynomial divisibility of height-one factorial
ratios}, Part I, in VladimirReshetnikov/ProveIt. Proposition
\texttt{prop:positive} supplies a smaller lcm multiplier, and Corollary
\texttt{cor:classificationAB} proves a stronger polynomial-divisor
classification. The original dated commit contains both statements.
The report explicitly has no Lean or Rocq formalization. This predecessor
was located after our initial release. The present contribution is
therefore an independent formalization of prior mathematics, not a
new mathematical resolution.

\begin{theorem}[Uniform positive-shift divisibility]
\label{thm:a211417-positive-shifts}
For every $k\in\{2,3,5\}$ and all integers $r,n\ge0$,
\[
 \prod_{\substack{1\le i\le r\\\gcd(i,k)=1}}(kn+i)
 \ \mid\ ((30r)!)^r A(n).
\]
In particular, Bala's positive-shift conjecture holds with the explicit
positive integer $C(k,r)=((30r)!)^r$.
\end{theorem}

This theorem independently formalizes the entire previously proved
positive-shift family, with arbitrarily many factors. The original
integrality of $A(n)$ is also a known result.  The
four fixed-factor assertions in the same OEIS comment have prior
kernel-checked proofs by Kenta Kitamura, and the separate negative-shift
$D(r)$ family has a reported paper-complete proof.  Neither is counted
as a new result here. The positive-shift family is likewise excluded
from the count of new resolutions. The theorem does not assert that the displayed
constant is minimal.

\subsection{Floor contributions and the large-divisor lemma}

For positive integers $m$, define
\[
 F(n,m)=\left\lfloor\frac{30n}{m}\right\rfloor
 +\left\lfloor\frac{n}{m}\right\rfloor
 -\left\lfloor\frac{15n}{m}\right\rfloor
 -\left\lfloor\frac{10n}{m}\right\rfloor
 -\left\lfloor\frac{6n}{m}\right\rfloor.
\]
Nested integer division, with $t=\lfloor30n/m\rfloor$, gives
\[
 F(n,m)=f(t),\qquad
 f(t)=t+\lfloor t/30\rfloor-\lfloor t/2\rfloor
                     -\lfloor t/3\rfloor-\lfloor t/5\rfloor.
\]
The function $f$ on nonnegative integers is periodic modulo 30, and
inspection of its 30 residues shows that $f(t)\in\{0,1\}$.
For clarity, the residues with value 1 are
\[
 \{1,2,3,4,5,7,8,9,11,13,14,17,19,23,29\}.
\]
Thus $F(n,m)\ge0$ for every $n,m$.  Legendre's formula then gives the
known integrality of $A(n)$ and the exact valuation identity
\begin{equation}
\label{eq:a211417-valuation}
 v_p(A(n))=\sum_{e\ge1}F(n,p^e)
\end{equation}
for every prime $p$.  All these sums are finite.

\begin{lemma}
\label{lem:a211417-large-divisor}
Suppose $k\in\{2,3,5\}$, $i\ge1$, $\gcd(i,k)=1$, and $m>30i$.
If $m\mid kn+i$, then $F(n,m)=1$.
\end{lemma}
\begin{proof}
Write $kn+i=mj$ and put $d=30/k$.  If $k\mid j$, then the same equation
implies $k\mid i$, contradicting $\gcd(i,k)=1$.  Therefore
$j\not\equiv0\pmod k$.
Multiplying the equation by $d$ gives
\[
 30n+di=mdj.
\]
Because $0<di<m$, we have
\[
 t=\left\lfloor\frac{30n}{m}\right\rfloor=dj-1.
\]
For $k=2$ the resulting residue of $t$ modulo 30 is 14.
For $k=3$ the possible residues are 9 and 19.
For $k=5$ they are 5, 11, 17, and 23.
Every one of these residues has $f(t)=1$, proving the claim.
\end{proof}

\subsection{Proof for the entire product}

\begin{proof}[Proof of Theorem~\ref{thm:a211417-positive-shifts}]
Fix $k,r,n$ as in the theorem and let
\[
 S=\{i\in\mathbb Z:1\le i\le r,\ \gcd(i,k)=1\},
 \qquad N_m=\#\{i\in S:m\mid kn+i\}.
\]
If $m>30r$, there is at most one $i\in S$ with $m\mid kn+i$.
Indeed, two distinct such indices would imply $m\mid i-j$, whereas
$0<|i-j|<m$.  If one such index exists, Lemma~\ref{lem:a211417-large-divisor}
applies because $i\le r$, and gives $F(n,m)=1$.  Consequently
\[
 N_m\le F(n,m)\qquad(m>30r).
\]
For $m\le30r$ we have $N_m\le|S|\le r$, while
$\lfloor30r/m\rfloor\ge1$ and $F(n,m)\ge0$.  These two cases combine as
\begin{equation}
\label{eq:a211417-uniform-contribution}
 N_m\le F(n,m)+r\left\lfloor\frac{30r}{m}\right\rfloor
 \qquad(m\ge1).
\end{equation}
For a fixed prime $p$, each factor $kn+i$ contributes one to its
valuation for each positive $e$ with $p^e\mid kn+i$.  Therefore
\[
 v_p\left(\prod_{i\in S}(kn+i)\right)=\sum_{e\ge1}N_{p^e}.
\]
This identity counts repeated prime factors across different selected
indices.  Summing~\eqref{eq:a211417-uniform-contribution} at $m=p^e$,
and using~\eqref{eq:a211417-valuation} and Legendre's formula again,
yields
\begin{align*}
 v_p\left(\prod_{i\in S}(kn+i)\right)
 &\le \sum_{e\ge1}F(n,p^e)
       +r\sum_{e\ge1}\left\lfloor\frac{30r}{p^e}\right\rfloor\\
 &=v_p(A(n))+r\,v_p((30r)!)\\
 &=v_p\bigl(((30r)!)^r A(n)\bigr).
\end{align*}
Since this holds for every prime, it proves the desired divisibility.
At $r=0$, the selected set is empty and $((30r)!)^r=1$, so the same
statement also includes the empty-product boundary.
\end{proof}

\subsection{Formalization and provenance}

The complete proof is formalized in \texttt{A211417/Proof.lean}, namespace
\texttt{ChebyshevShifts}, using Lean 4.19.0 and mathlib commit
\texttt{c44e0c8ee63ca166450922a373c7409c5d26b00b}.
The original numerator and denominator are defined with natural-number
factorials.  The theorem \texttt{factorial\_ratio\_identity} proves
that the quotient is exact; \texttt{sourceRatio\_eq\_a} then identifies
it with the literal rational factorial expression.
The endpoint \texttt{positive\_shift\_divisibility} proves the displayed
divisibility with the explicit constant, and
\texttt{positive\_shift\_rational\_integrality} constructs an integer
equal to the source's requested rational quotient.

The valuation proof is written with an explicit finite common cutoff
greater than $30n$, $30r$, and $kn+r$.  It controls the product directly,
including every multiplicity at every prime-power level.  The proof
contains no assumed conjecture, artificial finite bound, additional
axiom, or executable proof oracle.  The checked endpoints use only
\texttt{propext}, \texttt{Classical.choice}, and \texttt{Quot.sound}.

The source archive and bounded prior-resolution checks are supplied in
\texttt{A211417/source\_checks.md}.  The official OEIS bytes were checked
against the immutable export commit
\texttt{dff800c40a409482c9b734eecd1958986302e385}.  The initial searches missed the explicit predecessor. A broader
GitHub semantic code search later identified the September 19 paper
and its exact positive-family proposition. The package's
\texttt{PRIOR\_PROOF.md} records the original commit, exact statement
overlap, and chronology; \texttt{landau\_provenance.md} also explains
the classical multivariate Landau implication. The corrected status
is independent formalization of a previously proved assertion.

Primary links: \url{https://oeis.org/A211417};
\url{https://github.com/KitaKen1/oeis-a211417};
\url{https://github.com/google-deepmind/formal-conjectures/pull/5010};
\url{https://github.com/google-deepmind/formal-conjectures/issues/4923}.


Prior proof: OpenAI ChatGPT, \emph{Polynomial divisibility of height-one
factorial ratios}, Part I, September 19, 2026, unrefereed AI-assisted
research report. \href{https://github.com/VladimirReshetnikov/ProveIt/blob/a3fe9660ee6c3ab92752594c2099c3cd214012c9/docs/reports/congruences-and-valuations/factorial-ratio-polynomial-divisibility/article.tex}{Original dated source}, commit
\texttt{a3fe9660ee6c3ab92752594c2099c3cd214012c9}.
Classical background: Jonathan W. Bober, \emph{Factorial ratios,
hypergeometric series, and a family of step functions}, Theorem 3.1,
\url{https://arxiv.org/abs/0709.1977}.
